\begin{proposition}\label{prop:decic75-exclusion}
There is no elliptic curve over a number field of degree $10$
with a point of order $75$.
\end{proposition}
\begin{proof}
Put $X=X_1(75)$, of genus $145$. By \cite[Theorem~3.1]{DEHMZ},
$\operatorname{rank}J_X(\mathbb Q)=0$.
A hypothetical level structure gives a rational effective
trace divisor $D$ of degree $10$ on $X$, with noncuspidal support.

We use the quotient $\pi:X\to Y=X/\langle\delta_{26}\rangle$,
which parametrizes a point of order $25$ and a cyclic subgroup
of order $3$. The involution fixes the $25$-point, so has no
noncuspidal fixed point. Its action on the $112$ geometric cusps
is also free, including in characteristic $2$. Thus $\pi$ is
etale over $\mathbb Z_2$, and $Y$ has genus $73$ and good
reduction at $2$. Its rational Jacobian has rank zero by
pushforward. Write $\delta_u$ for the induced diamonds on $Y$.

The explicit reduced model is constructed over $X_1(25)$.
An exact modular unit gives a separable degree-$5$ presentation
$L/\mathbb F_2(t)$ of that genus-$12$ curve, with inverse
coordinates verified. For the universal $25$-point family put
$j=j_{25}$. In canonical ordinary coordinates, a cyclic
$3$-subgroup has common nonzero $x$-coordinate $w$ satisfying
\[
 w^4+w^3+1/j=0.
\]
This identity is verified directly from the division polynomial.
The quartic is irreducible, and its function field $M/L$ has
genus $73$. Thus this constructs the required quotient without
normalizing a genus-$145$ source model.

The full source Frobenius polynomial gives noncuspidal
closed-point counts
\[
 (0,0,0,0,0,0,20,0,0,120)
\]
in degrees $1$ through $10$. Consequently the noncuspidal part
of $\overline D$, if nonempty, is either one degree-$7$ place
with a cuspidal complement of degree $3$, or one degree-$10$
place. The cusps of degree at most $3$ are globally intact and
fixed by cyclotomic $7$.
Put $A_7=T_7-1-7\delta_7$, using the point-level geometric
isogeny correspondence. Eichler--Shimura and odd-prime torsion
injectivity give $A_7J_Y(\mathbb Q)=0$. At a standard rational
cusp $c_0$, the exact boundary formula gives $A_7c_0=0$.
It also gives $A_7C=0$ for the image of every possible degree-$3$
source cuspidal complement. Applying $A_7$ to
$[\pi_*D-10c_0]$ therefore forces the $A_7$-image of the
noncuspidal residual component to be principal.

There is one diamond orbit of degree-$7$ source points. The
retained exact computation on $Y$ gives
\[
 h^0(T_7P-P-7\delta_7P)=0,
\]
where $T_7P$ has four closed degree-$14$ terms, each with
multiplicity one. This computation is independent of the total
degree of $D$. Its boundary hypothesis remains valid for the
present degree-$3$ cuspidal complement, so excludes this branch.

In degree $10$ a fresh exhaustive census checks all $2053$
twists over $\mathbb F_{1024}$. Exactly $60$ have exponent
divisible by $75$. Their $j$-invariants form six Galois orbits
of degree $10$, and each curve has group $C_{1050}$. Indeed the
smaller invariant factor divides $1023$, and its square divides
$1050$, so it is $1$. There is consequently a unique cyclic
$75$-subgroup on each curve. Its generators, modulo sign, give
$20$ closed source points per $j$-orbit. The resulting $120$
points agree with the complete source Frobenius count.
They give six diamond orbits, and their images on $Y$ have
degree $10$, since their $j$-invariants already have that degree.

For one representative of each orbit we compute every actual
$7$-isogeny and carry both the $25$-point and the $3$-point.
The latter is used only through its cyclic subgroup on $Y$.
The kernel orbits have lengths $1,1,3,3$ over the source field;
the corresponding target places have degrees $10,10,30,30$,
with norm multiplicity $1$ in every case. Thus $T_7P$ has degree
$80$. The residue fields and ideal generators are checked
exactly. All six degree-zero divisors
\[
 B=T_7P-P-7\delta_7P
\]
have $h^0(B)=0$. A second native calculation rebuilds the
shared model and the $36$ prime ideals from their exact two
generators, checks the operator and norm multiplicities, and
repeats all six Riemann--Roch tests. Independent Sage calculations
also verify all $84$ actual moduli records, enumerate all $48$
geometric $7$-kernels, and check every actual level image by
Velu isogenies. In comparing geometric kernels with closed
places, the norm factors are applied individually; there is no
assumption that all kernel orbits have the same length.
This excludes noncuspidal reduction.

It remains to consider entirely cuspidal reduction. Here the
source and quotient cusp fields must be distinguished. The
source signatures over $\mathbb Q$ and $\mathbb F_2$ are
\[
 1^{20}2^{10}4^4 8^2 20^2,\qquad
 1^{20}2^{10}4^8 20^2.
\]
Some global octic cusps split into quartic places. On $Y$,
however, all cusp orbits are globally intact, with signature
$1^{20}4^4 10^2$. The complete primitive Tate labels verify
this descent as well as the following (degree,width) data:
\[
 (1,25)^{10},\quad(1,75)^{10},\quad
 (4,5)^2,\quad(4,15)^2,\quad(10,1),\quad(10,3).
\]
The source residue-degree factor is $2$ at the width-$25$
rational cusps and at both degree-$10$ cusps, and is $1$
elsewhere. Thus the width-$25$ multiplicities in
$\pi_*\overline D$ are even, and the degree-$10$ target cusps
cannot occur: their minimum source contribution has degree
$20$. The number of admissible images is
\[
 [t^{10}](1-t)^{-10}(1-t^2)^{-10}(1-t^4)^{-4}=871630.
\]

Let $\Omega=H^0(Y_{\mathbb F_2},\Omega^1)$ and let
$\mathcal C$ denote Cartier. The full-Jacobian endomorphism
$T_2-\delta_2$ has cotangent image
\[
 U=(\mathcal C-1)\Omega,\qquad \dim U=71.
\]
The Cartier matrix has rank $69$. The diamond $\delta_2$ has
order $10$, commutes with Cartier, and acts transitively on
each of the two sets of ten rational cusps. Its remaining
relevant cusp orbits have lengths $2,2$. These permutations are
independently matched with the complete Tate-label action.

The accepted common-basis packet contains genuine local
restriction kernels through weighted multiplicity $8$, together
with all diamond transports. A fresh complete binary enumeration
checks every admissible degree-$10$ divisor whose individual
weighted multiplicities are at most $8$: all $871520$
restriction matrices from $U$ have rank $10$. The omitted
$110$ configurations are $10c$ at a rational cusp, and $9c+r$
for distinct width-$75$ rational cusps $c,r$.

We compute these remaining ranks in a fresh single geometric
basis. The following direct ideal construction avoids a costly
cusp decomposition. Choose a finite rational cusp $p$ of $L$
with $v_p(j)=-25$, and let $I_B=\mathfrak p\mathcal O_M$.
The quartic has one unramified branch at $w=1$. The other
Newton slope is $-25/3$, so the remaining branch is totally
ramified of degree $3$, with residue field $\mathbb F_2$.
Denote the corresponding width-$25$ and width-$75$ cusps by
$c_{25},c_{75}$. Then
\[
 \pi_L^*p=c_{25}+3c_{75}.
\]
Globally the prime ideal of the first branch is obtained by
\[
 I_{25}=(I_B+(w-1)\mathcal O_M)\cap\mathcal O_M.
\]
The intersection removes unrelated poles of $w-1$. Native checks
verify primality, degree $1$, unramified restriction to $p$, and
$j$-width $25$. The integral ideal $I_B/I_{25}$ therefore
represents $3c_{75}$ directly.

In this fresh basis we verify the actual diamond pullback and
its commutation with Cartier. Exact differential spaces give
\[
 \begin{split}
 \operatorname{rank}(U|_{12c_{75}})&=12,\\
 \operatorname{rank}(U|_{10c_{25}})&=10,\\
 \operatorname{rank}(U|_{12c_{75}+3\delta_2^k c_{75}})&=15
       \qquad(1\leq k\leq9).
 \end{split}
\]
Surjectivity of restriction to an effective divisor implies
surjectivity on every effective subdivisor. These ranks imply
rank $10$ at $10c_{75}$ and at every
$9c_{75}+\delta_2^k c_{75}$. Diamond transitivity covers all
$110$ omitted configurations. Independent binary algebra checks
the fresh kernels, Cartier image, actual diamond matrix, and
all nine intersections. No matrices from different geometric
exports are combined.

Finally let $C$ be the rational cuspidal divisor on $Y$ lifting
$\pi_*\overline D$. The class $[\pi_*D-C]$ is rational torsion
and reduces to zero. The dyadic formal-group kernel lemma for
the full Jacobian shows that $T_2-\delta_2$ kills this class.
The rank-$10$ restriction tests above show that its Abel--Jacobi
composition on $Y^{(10)}$ is a formal immersion at $\overline C$.
It follows that $\pi_*D=C$, contradicting its noncuspidal
support. This uses the full Jacobian cotangent space and no
quotient-cotangent exactness assertion at $2$.
\end{proof}
